\documentclass{beamer}

\usepackage[utf8]{inputenc}

\input{../helper}
\graphicspath{{../img/}}


\title{Introdução à Computação \\ busca linear vs busca binária}
\author{Alexandre Rademaker}
\date{}

\begin{document}

\begin{frame}
 \maketitle
\end{frame}


\begin{frame}[fragile,allowframebreaks]{Busca linear}

  Queremos localizar um valor em um array. O computador só consegue
  examinar uma posição de memória por vez!

  Se nada pudermos assumir sobre a ordem dos valores no array, temos
  que olhar posição por posição. Buscando por 0.

  \begin{lstlisting}[basicstyle=\ttfamily]
  For each position from left to right
    If number is equal 0
      Return true
  Return false  
  \end{lstlisting}

  Busca linear é $O(n)$ e $\Theta(n)$. Por que?

  \framebreak

  \lstinputlisting[style=normalc,caption=src/numbers1.c]{src/numbers1.c}

  \framebreak

  \begin{lstlisting}[style=normalc,caption=src/names0.c]
    int main() {
      string names[] =
        {"Bill", "Charlie", "Fred",
         "George", "Ginny", "Percy", "Ron"};
  
       for (int i = 0; i < 7; i++) {
         if (names[i] == "Ron") {  // BUG
           printf("Found\n");
           return 0;
         }
       }
       printf("Not found\n");
       return 1;
     }
  \end{lstlisting}

  \framebreak

  \begin{lstlisting}[style=normalc,caption=src/names1.c]
    int main() {
      string names[] =
        {"Bill", "Charlie", "Fred",
         "George", "Ginny", "Percy", "Ron"};
      
       for (int i = 0; i < 7; i++) {
         if (strcmp(names[i], "Ron") == 0) {
           printf("Found\n");
           return 0;
         }
       }
       printf("Not found\n");
       return 1;
     }
  \end{lstlisting}

  \framebreak

  {\color{red} atenção} Usando o código abaixo teríamos o oposto do
  que queremos, pois valores diferentes de zero seriam considerados
  `true'. Vide \href{https://man7.org/linux/man-pages/man3/strcmp.3.html}{strcmp}.

  \begin{lstlisting}[style=normalc]
    if(strcmp(names[i], "Ron")) { ... }
  \end{lstlisting}

  Usando a negação funcionaria neste caso, mas seria um código pior,
  pois não verifica explicitamente o valor 0, que desejamos.

  \begin{lstlisting}[style=normalc]
    if(!strcmp(names[i], "Ron")) { ... }
  \end{lstlisting}

\end{frame}


\begin{frame}[fragile,allowframebreaks]{Busca binária}

  Queremos localizar um valor em um array. O computador só consegue
  examinar uma posição de memória por vez!

  Se pudermos assumir que os valores estão \textbf{ordenados}, podemos
  usar a busca binária. Buscando por N.

  \begin{lstlisting}[basicstyle=\ttfamily\scriptsize]
    If no position
     Return false
    If number in middle position equals N
     Return true
    Else if N < number in the middle position
     Search left half
    Else if N > number in the middle position
     Search right half
  \end{lstlisting}

  Busca binária é $O(\log n)$ e $\Omega(1)$. Por que? 

  \framebreak

  algumas considerações \ldots

  É muito mais eficiente que a busca linear, mas requer que o array
  esteja \textbf{ordenado}.

  Podemos ordenar o array antes\ldots Mas isso também irá demandar
  tempo. Se várias buscas forem ser executadas, a ordenação pode valer
  à pena (bancos de dados).

\end{frame}  


\begin{frame}[fragile,allowframebreaks]{Busca Binária Interativa em C}

  \begin{lstlisting}[style=tinyc,caption=src/numbers2.c]
  int binary_search(int a[], int n, int x) {
    int low = 0, high = n - 1;
    while (low <= high) {
      int k = (high + low) / 2;
      if      (a[k] == x) return k;
      else if (a[k] <  x) low  = k + 1;
      else                high = k - 1;
    }
    return -1;
  }

  int main(int argc, string argv[]) {
    int q, a[] = {-31, 0, 1, 2, 2, 4, 65, 83, 99, 782};
  
    if (argc > 1)  q = atoi(argv[1]);
    else           return 1;
    
    int n = sizeof(a)/sizeof(a[0]), i = binary_search(a, n, q);
    if (i >= 0) {
      printf("%d is at index %d.\n", q, i);
    }
    else {
      printf("%d is not found.\n", q);
    }
    return 0;
  }
  \end{lstlisting}

  \framebreak

  \begin{center}
    \begin{tabular}{|r|r|r|r|r|r|r|r|r|r|r|}
      \hline
      v & -31 & 0 & 1 & 2 & 2 & 4 & 65 & 83 & 99 & 782\\
      \hline
      p & 0 & 1 & 2 & 3 & 4 & 5 & 6 & 7 & 8 & 9\\
      \hline
    \end{tabular}
  \end{center}

  Buscando por 99, números são posições do array
  
  \begin{center}
    \begin{tabular}{rrrr}
      low & high & k & n\\ \hline
      0 & 9 & 4 & 10\\
      5 & 9 & 7 & 10\\
      8 & 9 & 8 & 10\\
    \end{tabular}
  \end{center}

\end{frame}


\begin{frame}[fragile,allowframebreaks]{Observações}

  Mas o código C mostrado é \textbf{interativo}, bem diferente do
  pseudo-código que é \textbf{recursivo}!

  A
  \href{https://man7.org/linux/man-pages/man0/stdlib.h.0p.html}{stdlib.h}
  onde achamos a função
  \href{https://man7.org/linux/man-pages/man3/atoi.3p.html}{atoi}.

  O operador \href{https://en.wikipedia.org/wiki/Sizeof}{sizeof} de C.

  Podemos \textbf{instrumentar} o código para entender seu funcionamento?

  \framebreak

  Podemos prever este tipo de erro? Como preparar o programa para um
  usuário menos \textbf{colaborativo}?

  \vspace{1cm}
  
  \begin{lstlisting}[basicstyle=\ttfamily]
    > ./numbers2 hello
    0 is at index 1.
  \end{lstlisting}

\end{frame}


\end{document}

%%% Local Variables:
%%% mode: latex
%%% TeX-master: t
%%% End:
